\section{Introducción}

La especificación de requisitos permite definir qué debe hacer un sistema y las
condiciones que debe cumplir. Además de delimitar el alcance del proyecto, sirve
como referencia para comprobar que las funcionalidades previstas se han
implementado.

En este proyecto, la especificación se organiza a partir de los
\textbf{objetivos generales}, el \textbf{catálogo de requisitos}, los
\textbf{casos de uso} y la \textbf{trazabilidad} entre estos elementos.

La especificación toma como referencia la norma
\textbf{\acs{ISO}/\acs{IEC}/\acs{IEEE}
29148:2018}~\cite{iso_29148_2018}, que define los procesos y contenidos asociados
a la ingeniería de requisitos. Los requisitos de calidad se organizan a partir
del modelo de calidad de producto propuesto por \textbf{\acs{ISO}/\acs{IEC}
25010:2023}~\cite{iso_25010_2023}. Para hacer verificables estos requisitos, se
utiliza el marco definido por \textbf{\acs{ISO}/\acs{IEC}
25030:2019}~\cite{iso_25030_2019}. Las tres normas en conjunto se han aplicado
de forma adaptada al alcance académico y técnico del proyecto.

Como apoyo a la redacción, también se ha utilizado la lista de comprobación
publicada por la NASA, que recoge criterios prácticos para que cada requisito
sea claro, sencillo y verificable~\cite{nasa_good_requirement}.

\newcommand{\idObjetivo}[1]{%
    \hypertarget{obj:#1}{#1}%
}

\newcommand{\refObjetivo}[1]{%
    \hyperlink{obj:#1}{#1}%
}

\section{Objetivos generales}

El proyecto persigue los siguientes objetivos generales:

\begin{itemize}
    \item \textbf{\idObjetivo{OG-01}}: Desarrollar una aplicación web que facilite el
    aprendizaje y la consulta de las reglas de juegos de mesa.

    \item \textbf{\idObjetivo{OG-02}}: Procesar manuales en distintos formatos para extraer
    su contenido y poder utilizarlo dentro de la aplicación.

    \item \textbf{\idObjetivo{OG-03}}: Generar explicaciones y responder preguntas sobre un
    juego a partir de la información contenida en sus manuales.

    \item \textbf{\idObjetivo{OG-04}}: Almacenar de forma organizada la información de los
    juegos, los manuales y las conversaciones de cada usuario.

    \item \textbf{\idObjetivo{OG-05}}: Garantizar la seguridad de la
    aplicación y proteger la información de los usuarios.
\end{itemize}

El recomendador de juegos queda fuera del alcance del proyecto para ajustar
el desarrollo al tiempo disponible.

\section{Catálogo de requisitos}

\newcommand{\idRequisito}[1]{%
    \hypertarget{req:#1}{#1}%
}

\newcommand{\refRequisito}[1]{%
    \hyperlink{req:#1}{#1}%
}

\newcommand{\subrequisito}[3]{%
    \par\smallskip
    \textbf{\idRequisito{#1}. #2}: #3%
}

\newcommand{\fichaRequisito}[9]{%
    \begin{table}[H]
    \centering
    \begin{tblr}{
        width=\linewidth,
        colspec={Q[l,wd=0.26\linewidth] X[l]},
        rows={valign=t},
        row{1}={bg=gray!18,font=\bfseries},
        column{1}={font=\bfseries},
        hline{1}={0.8pt},
        hline{2}={0.5pt},
        hline{Z}={0.8pt},
        leftsep=4pt,
        rightsep=4pt,
        rowsep=2pt
    }
        \idRequisito{#1} & #2 \\
        Objetivo general & #3 \\
        #4 & #5 \\
        Descripción & #6 \\
        Prioridad & #7 \\
        Relacionado con & #8 \\
        #9%
    \end{tblr}
    \caption{\textbf{#1}. #2.}
    \label{tab:req:#1}
    \end{table}
}

\newcommand{\requisitoFuncional}[8]{%
    \fichaRequisito{#1}{#2}{#3}{Requisito padre}{\refRequisito{#4}}%
    {#5}{#6}{#7}{Verificación & #8 \\}%
}

\newcommand{\requisitoPadre}[8]{%
    \fichaRequisito{#1}{#2}{#3}{Requisitos hijos}{#4}{#5}{#6}{#7}%
    {Verificación & #8 \\}%
}



Los requisitos funcionales describen \textbf{operaciones}, \textbf{resultados} y \textbf{cambios de
estado}, por otro lado, los no funcionales establecen \textbf{condiciones de calidad}.

Cada requisito está relacionado con los objetivos que ayuda a cumplir y con
otros requisitos cuando el funcionamiento de uno afecta al otro. Estas
relaciones permiten mantener una \textbf{trazabilidad} de los requisitos del
proyecto~\cite{nasa_good_requirement}.


\newcommand{\requisitoNoFuncional}[7]{%
    \fichaRequisito{#1}{#2}{#3}{Alcance}{#4}{#5}{#6}{#7}{}%
}

\setsecnumdepth{subsection}
\settocdepth{subsection}
\subsection{Requisitos funcionales}

A partir de los objetivos generales, se han extraído una serie de \textbf{requisitos
funcionales} específicos que establecen cuál debe ser el comportamiento del
sistema.

Para estructurarlos, se han definido cinco requisitos principales, de los que parte el resto:


\requisitoPadre
{RF-1}
{Gestión de cuentas}
{\refObjetivo{OG-04} y \refObjetivo{OG-05}}
{\refRequisito{RF-1.1}, \refRequisito{RF-1.2}, \refRequisito{RF-1.3},
\refRequisito{RF-1.4}, \refRequisito{RF-1.5}, \refRequisito{RF-1.6},
\refRequisito{RF-1.7}, \refRequisito{RF-1.8} y \refRequisito{RF-1.9}}
{El sistema debe permitir gestionar el ciclo de vida completo de una cuenta.}
{Alta}
{\refRequisito{RF-2}, \refRequisito{RF-3} y \refRequisito{RF-4}}
{Completar el ciclo de una cuenta, desde su registro hasta su eliminación, y
comprobar durante el proceso las operaciones de acceso, verificación,
recuperación, edición y consulta de información asociada.}

\requisitoPadre
{RF-2}
{Gestión de juegos}
{\refObjetivo{OG-01} y \refObjetivo{OG-04}}
{\refRequisito{RF-2.1}, \refRequisito{RF-2.2}, \refRequisito{RF-2.3},
\refRequisito{RF-2.4}, \refRequisito{RF-2.5} y \refRequisito{RF-2.6}}
{El sistema debe permitir localizar y registrar juegos de mesa, consultar su información
y organizar la actividad del usuario asociada a ellos.}
{Alta}
{\refRequisito{RF-1}, \refRequisito{RF-3} y \refRequisito{RF-4}}
{Buscar y registrar un juego, seguirlo y dejar de seguirlo, consultar su ficha y
la biblioteca, y crear, actualizar y retirar una valoración.}

\requisitoPadre
{RF-3}
{Procesamiento de manuales}
{\refObjetivo{OG-02} y \refObjetivo{OG-04}}
{\refRequisito{RF-3.1}, \refRequisito{RF-3.2}, \refRequisito{RF-3.3},
\refRequisito{RF-3.4}, \refRequisito{RF-3.5}, \refRequisito{RF-3.6},
\refRequisito{RF-3.7}, \refRequisito{RF-3.8}, \refRequisito{RF-3.9} y
\refRequisito{RF-3.10}}
{El sistema debe permitir añadir, procesar, consultar y gestionar manuales
asociados a juegos de mesa.}
{Alta}
{\refRequisito{RF-1}, \refRequisito{RF-2} y \refRequisito{RF-4}}
{Subir manuales de juegos de mesa y comprobar que las subidas inválidas o duplicadas se
rechazan. Después, procesarlos, consultar su estado y contenido, corregir y
reprocesar sus páginas y eliminarlos.}

\requisitoPadre
{RF-4}
{Consulta de reglas}
{\refObjetivo{OG-03} y \refObjetivo{OG-04}}
{\refRequisito{RF-4.1}, \refRequisito{RF-4.2}, \refRequisito{RF-4.3},
\refRequisito{RF-4.4}, \refRequisito{RF-4.5}, \refRequisito{RF-4.6},
\refRequisito{RF-4.7} y \refRequisito{RF-4.8}}
{El sistema debe permitir obtener explicaciones y mantener conversaciones
sobre las reglas de un juego de mesa utilizando el contenido de los manuales
disponibles.}
{Alta}
{\refRequisito{RF-1}, \refRequisito{RF-2} y \refRequisito{RF-3}}
{Generar la explicación de un juego y mantener una conversación sobre sus
reglas con un modelo \ac{LLM}.}


\requisitoPadre
{RF-5}
{Gestión de preferencias}
{\refObjetivo{OG-01}}
{\refRequisito{RF-5.1} y \refRequisito{RF-5.2}}
{El sistema debe permitir adaptar el idioma y la apariencia de la interfaz a las
preferencias del usuario.}
{Baja}
{Sin relación directa}
{Cambiar el idioma y la apariencia de la interfaz, recargar la aplicación
y comprobar que se conservan las preferencias seleccionadas.}


\subsubsection{RF-1. Gestión de cuentas}

La gestión de una cuenta cubre todo su ciclo, desde el registro hasta la
eliminación, e incluye el acceso, la verificación, la recuperación y la edición
de sus datos entre otros detalles.

\requisitoFuncional
{RF-1.1}
{Registrar una cuenta nueva}
{\refObjetivo{OG-04} y \refObjetivo{OG-05}}
{RF-1}
{El sistema debe permitir crear una cuenta mediante un nombre de usuario, un
correo electrónico y una contraseña. Tras completar el registro, la sesión debe
quedar iniciada.}
{Alta}
{\refRequisito{RF-1.2} y \refRequisito{RF-1.4}}
{Completar un registro con datos válidos y comprobar que se crea la cuenta y se
inicia la sesión.}

\requisitoFuncional
{RF-1.2}
{Iniciar sesión}
{\refObjetivo{OG-04} y \refObjetivo{OG-05}}
{RF-1}
{El sistema debe permitir iniciar sesión utilizando el correo
electrónico, o bien, el nombre de usuario junto con la contraseña.}
{Alta}
{\refRequisito{RF-1.1}, \refRequisito{RF-1.3}, \refRequisito{RF-1.5},
\refRequisito{RF-1.6}, \refRequisito{RF-1.7} y \refRequisito{RF-1.8}}
{Acceder a una misma cuenta mediante ambos identificadores y rechazar unas
credenciales que no sean válidas.}

\requisitoFuncional
{RF-1.3}
{Cerrar sesión}
{\refObjetivo{OG-04} y \refObjetivo{OG-05}}
{RF-1}
{El sistema debe permitir cerrar la sesión y, una vez cerrada, no debe
poder utilizarse de nuevo sin autenticarse.}
{Alta}
{\refRequisito{RF-1.2}}
{Cerrar una sesión activa y comprobar que no puede volver a utilizarse para
acceder a una ruta protegida.}

\requisitoFuncional
{RF-1.4}
{Verificar el correo electrónico}
{\refObjetivo{OG-04} y \refObjetivo{OG-05}}
{RF-1}
{El sistema debe permitir verificar el correo mediante un enlace de un solo uso
y solicitar un nuevo enlace cuando sea necesario. La falta de verificación no
debe impedir el acceso al resto de la aplicación.}
{Alta}
{\refRequisito{RF-1.1} y \refRequisito{RF-1.6}}
{Verificar una cuenta, rechazar la reutilización del enlace y comprobar que una
cuenta pendiente de verificación puede seguir utilizando la aplicación.}

\requisitoFuncional
{RF-1.5}
{Restablecer la contraseña}
{\refObjetivo{OG-04} y \refObjetivo{OG-05}}
{RF-1}
{El sistema debe permitir restablecer la contraseña mediante un enlace de un
solo uso enviado por correo electrónico. Al completar el proceso, debe revocar
las sesiones que continuasen activas.}
{Alta}
{\refRequisito{RF-1.2} y \refRequisito{RF-1.6}}
{Restablecer la contraseña, comprobar que el enlace no puede reutilizarse y que
las sesiones anteriores dejan de ser válidas.}

\requisitoFuncional
{RF-1.6}
{Editar el perfil}
{\refObjetivo{OG-04}}
{RF-1}
{El sistema debe permitir actualizar de forma independiente el nombre de
usuario, el correo electrónico y la imagen de perfil. Un cambio de correo debe
dejar pendiente su verificación y generar un nuevo enlace.}
{Media}
{\refRequisito{RF-1.2}, \refRequisito{RF-1.4} y \refRequisito{RF-1.5}}
{Modificar por separado cada dato del perfil y comprobar el estado de
verificación después de cambiar el correo electrónico.}

\requisitoFuncional
{RF-1.7}
{Cambiar la contraseña}
{\refObjetivo{OG-04} y \refObjetivo{OG-05}}
{RF-1}
{El sistema debe permitir cambiar la contraseña después de comprobar la actual
y validar la nueva. La sesión utilizada para realizar el cambio debe ser la única
que sigue activa tras el cambio.}
{Alta}
{\refRequisito{RF-1.2}}
{Cambiar la contraseña desde una sesión con otras sesiones abiertas y comprobar
que solo permanece activa la cuenta desde la que se ha hecho el cambio.}

\requisitoFuncional
{RF-1.8}
{Eliminar una cuenta}
{\refObjetivo{OG-04} y \refObjetivo{OG-05}}
{RF-1}
{El sistema debe permitir eliminar una cuenta tras realizar una confirmación
explícita. Una vez eliminada, se deben revocar sus sesiones y hacer inaccesibles
la cuenta, los manuales, las conversaciones y las valoraciones asociadas. }
{Alta}
{\refRequisito{RF-1.2}, \refRequisito{RF-2.6}, \refRequisito{RF-3.9} y
\refRequisito{RF-4.2}}
{Eliminar una cuenta con información asociada y comprobar que deja de ser
accesible, que sus sesiones se cierran y que se hace la limpieza de sus
recursos.}

\requisitoFuncional
{RF-1.9}
{Consultar las estadísticas de actividad}
{\refObjetivo{OG-04}}
{RF-1}
{El sistema debe mostrar al usuario el número de juegos en los que ha tenido
actividad, los manuales que ha subido y las conversaciones que ha creado.}
{Media}
{\refRequisito{RF-2.6}, \refRequisito{RF-3.1}, \refRequisito{RF-3.9} y
\refRequisito{RF-4.2}}
{Generar actividad de todo tipo y comprobar que los valores mostrados
coinciden con la información almacenada.}

\subsubsection{RF-2. Gestión de juegos}

Los juegos pueden buscarse en el catálogo o añadirse manualmente. Una vez
registrados, pueden seguirse, valorarse y consultarse desde la biblioteca
personal.

\requisitoFuncional
{RF-2.1}
{Buscar juegos}
{\refObjetivo{OG-01}}
{RF-2}
{El sistema debe permitir buscar juegos por nombre.}
{Alta}
{\refRequisito{RF-2.2} y \refRequisito{RF-3.1}}
{Buscar primero un juego que ya haya tenido actividad y otro que no para
comprobar que ambos se muestran.}

\requisitoFuncional
{RF-2.2}
{Añadir un juego manualmente}
{\refObjetivo{OG-04}}
{RF-2}
{El sistema debe permitir añadir mediante su nombre un juego que no aparezca
entre los resultados de búsqueda.}
{Media}
{\refRequisito{RF-2.1} y \refRequisito{RF-3.1}}
{Buscar un juego que no esté disponible, añadirlo manualmente y comprobar que
puede seleccionarse para subir un manual.}

\requisitoFuncional
{RF-2.3}
{Permitir seguir juegos}
{\refObjetivo{OG-04}}
{RF-2}
{El sistema debe registrar los juegos seguidos de un usuario
\subrequisito{RF-2.3.1}{Seguir o dejar de seguir un juego}{El usuario debe poder seguir y
dejar de seguir un juego.}
\subrequisito{RF-2.3.2}{Seguir juegos automáticamente}{La primera actividad sobre un
juego debe seguirlo automáticamente cuando no exista una decisión explícita
anterior.}}
{Media}
{\refRequisito{RF-2.4}, \refRequisito{RF-2.5}, \refRequisito{RF-2.6},
\refRequisito{RF-3.1} y \refRequisito{RF-4.3}}
{Seguir y dejar de seguir un juego de forma explícita. Después, generar una
primera actividad con y sin una decisión anterior y comprobar el resultado.}

\requisitoFuncional
{RF-2.4}
{Consultar la biblioteca personal}
{\refObjetivo{OG-04}}
{RF-2}
{El sistema debe recoger la actividad del usuario relacionada con juegos y
manuales.
\subrequisito{RF-2.4.1}{Juegos seguidos}{Debe listar los juegos a los que
sigue el usuario.}
\subrequisito{RF-2.4.2}{Manuales propios}{Debe listar los manuales del usuario subidos por el usuario y permitir filtrarlos dentro de la biblioteca.}}
{Media}
{\refRequisito{RF-2.3}, \refRequisito{RF-3.1} y \refRequisito{RF-3.9}}
{Seguir varios juegos y subir varios manuales. Después, comprobar el orden de
los juegos y el filtrado de los manuales.}

\requisitoFuncional
{RF-2.5}
{Consultar el perfil de un juego}
{\refObjetivo{OG-01}}
{RF-2}
{El sistema debe mostrar los datos del juego, los manuales accesibles, la
valoración propia y el número de conversaciones del usuario, e indicar si sigue el juego.}
{Alta}
{\refRequisito{RF-2.3}, \refRequisito{RF-2.6}, \refRequisito{RF-3.1},
\refRequisito{RF-3.4}, \refRequisito{RF-3.9} y \refRequisito{RF-4.2}}
{Consultar un juego con información de cada tipo y comprobar que la ficha
solo incluye los manuales accesibles para la cuenta autenticada.}

\requisitoFuncional
{RF-2.6}
{Valorar un juego}
{\refObjetivo{OG-04}}
{RF-2}
{El sistema debe permitir registrar y actualizar una puntuación de entre
uno y cinco a un juego, acompañarla de una nota opcional y permitir
borrar valoraciones previas.}
{Media}
{\refRequisito{RF-1.8}, \refRequisito{RF-1.9}, \refRequisito{RF-2.3} y
\refRequisito{RF-2.5}}
{Crear, modificar y borrar una valoración.}

\subsubsection{RF-3. Procesamiento de manuales}

Los manuales se validan al subirlos y se procesan antes de poder ser
consultados.
También pueden corregirse manualmente errores, reprocesarse y eliminarse.

\requisitoFuncional
{RF-3.1}
{Subir un manual}
{\refObjetivo{OG-02}}
{RF-3}
{El sistema debe permitir subir un manual asociado a un juego mediante
imágenes o un documento en formato \acs{PDF}\acused{PDF}, y elegir si será
privado o compartido con la comunidad. Si se comparte, debe permitir elegir
si se muestra el nombre del usuario, de acuerdo con \refRequisito{RF-3.10}.
Antes de enviar varias imágenes, el usuario debe poder revisarlas, ordenarlas y retirar las que no sean útiles.}
{Media}
{\refRequisito{RF-1.9}, \refRequisito{RF-2.1}, \refRequisito{RF-2.2},
\refRequisito{RF-2.3}, \refRequisito{RF-2.4}, \refRequisito{RF-2.5},
\refRequisito{RF-3.2}, \refRequisito{RF-3.3}, \refRequisito{RF-3.4},
\refRequisito{RF-3.6.1}, \refRequisito{RF-3.8}, \refRequisito{RF-3.9} y
\refRequisito{RF-3.10}}
{Subir un manual de cada tipo y comprobar que se
conserva el orden de las páginas, la visibilidad y la opción elegida para
mostrar el nombre.}

\requisitoFuncional
{RF-3.2}
{Validar la subida de un manual}
{\refObjetivo{OG-02}}
{RF-3}
{El sistema debe comprobar en el servidor que los formatos admitidos, el tamaño de
cada archivo y del conjunto de la petición, y el número máximo de páginas
están dentro de los valores esperados.}
{Media}
{\refRequisito{RF-3.1} y \refRequisito{RF-3.4}}
{Subir manuales con valores correctos y otros que superen cada límite,
comprobando que estos últimas se rechazan antes de iniciar el procesamiento.}

\requisitoFuncional
{RF-3.3}
{Detectar manuales duplicados}
{\refObjetivo{OG-02} y \refObjetivo{OG-04}}
{RF-3}
{El sistema debe rechazar un manual que ya está activo y lo ha subido
el mismo usuario.}
{Media}
{\refRequisito{RF-3.1} y \refRequisito{RF-3.9}}
{Subir un manual dos veces como el mismo usuario y juego y comprobar que  el segundo se rechaza.}



\requisitoFuncional
{RF-3.4}
{Procesar un manual}
{\refObjetivo{OG-02}}
{RF-3}
{El sistema debe procesar y extraer el contenido de cada manual en segundo
plano para dejarlo preparado para las consultas al \ac{LLM}.
\subrequisito{RF-3.4.1}{Preprocesar las imágenes}{Cuando sea necesario reconocer
texto mediante \ac{OCR}, debe aplicar a las imágenes las técnicas de preprocesado adecuadas según el motor utilizado.}
\subrequisito{RF-3.4.2}{Extraer el texto}{Debe utilizar el texto del documento
cuando sea aprovechable. En caso contrario, debe extraerlo de las imágenes
preprocesadas mediante \ac{OCR}.}
\subrequisito{RF-3.4.3}{Posprocesar el texto}{Tras la extracción de texto,
debe aplicar las distintas técnicas de postprocesado para mejorar el
material.}
\subrequisito{RF-3.4.4}{Indexar el contenido}{Debe dividir el texto resultante en
fragmentos e indexarlos en una base de datos vectorial~\cite{microsoft_vector_database}
para recuperarlos durante las consultas.}
\subrequisito{RF-3.4.5}{Reutilizar páginas ya procesadas}{Debe reutilizar el resultado de una
página idéntica ya procesada, sin repetir su extracción ni su posprocesado.}}
{Media}
{\refRequisito{RF-2.5}, \refRequisito{RF-3.1}, \refRequisito{RF-3.2},
\refRequisito{RF-3.5}, \refRequisito{RF-3.6}, \refRequisito{RF-3.7},
\refRequisito{RF-3.8}, \refRequisito{RF-3.9}, \refRequisito{RF-4.1} y
\refRequisito{RF-4.3}}
{Procesar un documento con texto, una imagen que requiera reconocimiento y una
página repetida. Después, comprobar el texto obtenido y su indexación.}


\requisitoFuncional
{RF-3.5}
{Consultar el estado del procesamiento}
{\refObjetivo{OG-02}}
{RF-3}
{El sistema debe permitir al usuario propietario consultar el estado guardado del
manual y de cada una de sus páginas.}
{Media}
{\refRequisito{RF-3.4}}
{Iniciar un procesamiento y comprobar sus diferentes transiciones de estado
tanto para el manual como para sus páginas.}

\requisitoFuncional
{RF-3.6}
{Consultar un manual}
{\refObjetivo{OG-02} y \refObjetivo{OG-04}}
{RF-3}
{El sistema debe permitir consultar por páginas el contenido de un manual
propio o compartido que esté activo. En los manuales ajenos, la consulta
debe limitarse a la lectura.
\subrequisito{RF-3.6.1}{Navegar por las páginas}{Debe mostrar el texto extraído
o editado y, cuando esté disponible, la imagen subida o generada de cada página.}
\subrequisito{RF-3.6.2}{Buscar en el manual}{Debe buscar texto dentro del manual y
permitir desplazarse entre las coincidencias.}
\subrequisito{RF-3.6.3}{Permitir obtener información sobre el procesamiento}{Debe mostrar al propietario la información
disponible sobre el origen y la confianza del texto, así como la reutilización
de páginas.}}
{Alta}
{\refRequisito{RF-3.1}, \refRequisito{RF-3.4}, \refRequisito{RF-3.7},
\refRequisito{RF-3.9} y \refRequisito{RF-4.6}}
{Consultar un manual con varias páginas, buscar un fragmento de texto presente y comprobar
la navegación, el texto, la imagen disponible y la información técnica de cada
página.}

\requisitoFuncional
{RF-3.7}
{Corregir el texto de una página manualmente}
{\refObjetivo{OG-02}}
{RF-3}
{El sistema debe permitir corregir el texto de una página cuando el manual sea
propio, privado y no se esté procesando. Tras guardar el cambio, debe encolar
la actualización del contenido indexado.}
{Alta}
{\refRequisito{RF-3.4.4}, \refRequisito{RF-3.6.1}, \refRequisito{RF-4.1} y
\refRequisito{RF-4.3}}
{Editar una página que cumpla las condiciones y comprobar que se reindexa correctamente.}

\requisitoFuncional
{RF-3.8}
{Reprocesar un manual}
{\refObjetivo{OG-02}}
{RF-3}
{El sistema debe permitir repetir el procesamiento a partir de las fuentes que
ya se encuentran almacenadas.
\subrequisito{RF-3.8.1}{Reprocesar el manual completo}{Debe volver a procesar
todas sus páginas.}
\subrequisito{RF-3.8.2}{Reprocesar una página}{Debe procesar únicamente la página seleccionada.}}
{Media}
{\refRequisito{RF-3.1}, \refRequisito{RF-3.4}, \refRequisito{RF-3.9},
\refRequisito{RF-4.1} y \refRequisito{RF-4.3}}
{Reprocesar primero una página y después el manual completo, comprobando las
páginas afectadas subirlas de nuevo.}

\requisitoFuncional
{RF-3.9}
{Eliminar un manual}
{\refObjetivo{OG-04}}
{RF-3}
{El sistema debe permitir aplicar un borrado lógico a un manual propio y retirar
sus páginas del acceso del usuario. Después debe encolar la limpieza de los
ficheros y del contenido indexado asociado.}
{Alta}
{\refRequisito{RF-1.8}, \refRequisito{RF-1.9}, \refRequisito{RF-2.4},
\refRequisito{RF-2.5}, \refRequisito{RF-3.1}, \refRequisito{RF-3.3},
\refRequisito{RF-3.4.4}, \refRequisito{RF-3.6}, \refRequisito{RF-3.8},
\refRequisito{RF-4.1}, \refRequisito{RF-4.3}, \refRequisito{RF-4.6} y
\refRequisito{RF-4.8}}
{Eliminar un manual procesado y comprobar que deja de ser accesible y que se
solicita la limpieza de sus recursos derivados.}

\requisitoFuncional
{RF-3.10}
{Mostrar u ocultar el nombre al compartir un manual}
{\refObjetivo{OG-04} y \refObjetivo{OG-05}}
{RF-3}
{El sistema debe permitir al propietario elegir si su nombre aparece junto
a un manual compartido y en las fuentes de las respuestas que lo utilizan.
La opción por defecto debe ocultarlo y debe poder cambiarse después de la
subida, también para las fuentes de respuestas anteriores. Esta elección
no debe modificar la visibilidad del manual ni sus permisos.}
{Media}
{\refRequisito{RF-3.1}, \refRequisito{RF-3.6} y \refRequisito{RF-4.6}}
{Subir un manual compartido sin mostrar el nombre y comprobar su presentación
desde otra cuenta. Activar y desactivar la opción, verificando el cambio
en el manual y en una respuesta anterior. Rechazar el cambio desde una
cuenta ajena y mantener oculto el nombre en los manuales privados.}

\subsubsection{RF-4. Consulta de reglas}

Las explicaciones y respuestas del modelo \ac{LLM} se generan a partir de los manuales disponibles.
Las conversaciones conservan el historial y las fuentes utilizadas.

\requisitoFuncional
{RF-4.1}
{Consultar la explicación de un juego}
{\refObjetivo{OG-03}}
{RF-4}
{El sistema debe generar en segundo plano una explicación general a partir de
los manuales accesibles y conservarla mientras sus fuentes no cambien.}
{Alta}
{\refRequisito{RF-3.4.4}, \refRequisito{RF-3.7}, \refRequisito{RF-3.8} y
\refRequisito{RF-3.9}}
{Solicitar una explicación inexistente, comprobar su generación y reutilizarla
en una consulta posterior. Después, modificar las fuentes y comprobar que se
genera una nueva versión.}

\requisitoFuncional
{RF-4.2}
{Gestionar conversaciones}
{\refObjetivo{OG-03} y \refObjetivo{OG-04}}
{RF-4}
{El sistema debe permitir crear conversaciones persistentes con el modelo
\ac{LLM} asociadas a cada juego.
\subrequisito{RF-4.2.1}{Crear y listar}{Debe permitir crear una conversación y
listar las conversaciones propias del juego.}
\subrequisito{RF-4.2.2}{Consultar el historial}{Debe conservar y mostrar los
mensajes de una conversación.}
\subrequisito{RF-4.2.3}{Renombrar}{Debe permitir modificar el título de una
conversación propia.}
\subrequisito{RF-4.2.4}{Eliminar}{Debe permitir eliminar una conversación
propia.}}
{Alta}
{\refRequisito{RF-1.8}, \refRequisito{RF-1.9}, \refRequisito{RF-2.5},
\refRequisito{RF-4.3}, \refRequisito{RF-4.4}, \refRequisito{RF-4.7} y
\refRequisito{RF-4.8}}
{Crear una conversación, mandar mensajes, consultar su historial,
renombrarla y eliminarla.}

\requisitoFuncional
{RF-4.3}
{Generar una respuesta}
{\refObjetivo{OG-03}}
{RF-4}
{El sistema debe generar en segundo plano las respuestas a los mensajes de una
conversación y mostrar si se encuentran pendientes, completadas o fallidas.}
{Alta}
{\refRequisito{RF-2.3}, \refRequisito{RF-3.4.4}, \refRequisito{RF-3.7},
\refRequisito{RF-3.8}, \refRequisito{RF-3.9}, \refRequisito{RF-4.2},
\refRequisito{RF-4.4}, \refRequisito{RF-4.5}, \refRequisito{RF-4.6},
\refRequisito{RF-4.7} y \refRequisito{RF-4.8}}
{Enviar una pregunta y comprobar que termina con una respuesta o con un error
explícito, sin permanecer indefinidamente en estado pendiente.}

\requisitoFuncional
{RF-4.4}
{Nombrar una conversación}
{\refObjetivo{OG-03}}
{RF-4}
{El sistema debe asignar automáticamente un título a partir del primer mensaje
de una conversación.}
{Alta}
{\refRequisito{RF-4.2.3} y \refRequisito{RF-4.3}}
{Crear una conversación, enviar su primer mensaje y comprobar que el título
inicial se sustituye por uno generado a partir de dicho mensaje.}

\requisitoFuncional
{RF-4.5}
{Seleccionar el idioma de la respuesta}
{\refObjetivo{OG-03}}
{RF-4}
{El sistema debe responder a los mensajes de una conversación en español o en
inglés según el idioma detectado en el turno.}
{Alta}
{\refRequisito{RF-4.3}}
{Enviar preguntas equivalentes en ambos idiomas y comprobar el idioma de las
respuestas generadas.}

\requisitoFuncional
{RF-4.6}
{Consultar las fuentes de una respuesta}
{\refObjetivo{OG-03}}
{RF-4}
{El sistema debe indicar las páginas y el manual en el que se han basado
las respuestas dadas a una pregunta. Cuando el manual continúe disponible
y el usuario tenga permiso para leerlo, debe permitir abrir directamente
la página citada. El nombre de quien lo subió debe mostrarse según la
elección recogida en \refRequisito{RF-3.10}.}
{Alta}
{\refRequisito{RF-3.6.1}, \refRequisito{RF-3.9}, \refRequisito{RF-3.10} y \refRequisito{RF-4.3}}
{Generar una respuesta con fuentes procedentes de varias páginas y manuales, y
comprobar que cada referencia se muestra una sola vez y abre la página correcta
cuando se dispone de acceso.}

\requisitoFuncional
{RF-4.7}
{Reutilizar contexto de una conversación}
{\refObjetivo{OG-03}}
{RF-4}
{El sistema debe proveer del contexto de la conversación al \ac{LLM} tras
recibir un nuevo mensaje.}
{Alta}
{\refRequisito{RF-4.2.2} y \refRequisito{RF-4.3}}
{Realizar una pregunta en una conversación previamente iniciada que omita parte del contexto y comprobar
que se interpreta utilizando los mensajes anteriores.}

\requisitoFuncional
{RF-4.8}
{Conservar conversaciones sin fuentes disponibles}
{\refObjetivo{OG-03} y \refObjetivo{OG-04}}
{RF-4}
{El sistema debe mantener accesible el historial de una conversación aunque las fuentes que se hayan usado se retiren. En ese caso, no
debe admitir nuevos mensajes.}
{Alta}
{\refRequisito{RF-3.9}, \refRequisito{RF-4.2.2} y \refRequisito{RF-4.3}}
{Eliminar los manuales vinculados a una conversación existente y comprobar que
su historial continúa visible en modo de solo lectura.}

\subsubsection{RF-5. Gestión de preferencias}

Las preferencias de \textit{Manualito} permiten elegir entre español e
inglés y cambiar la apariencia de la interfaz.


\requisitoFuncional
{RF-5.1}
{Seleccionar el idioma de la interfaz}
{\refObjetivo{OG-01}}
{RF-5}
{El sistema debe ofrecer la interfaz en español e inglés y permitir cambiar el
idioma desde todas las partes de la aplicación. La selección debe
conservarse en el dispositivo.}
{Media}
{Sin relación directa}
{Cambiar el idioma antes y después de iniciar sesión, recargar la aplicación y
comprobar que la interfaz mantiene la selección.}

\requisitoFuncional
{RF-5.2}
{Configurar la apariencia}
{\refObjetivo{OG-01}}
{RF-5}
{El sistema debe permitir elegir entre un tema claro, oscuro o automático y
permitir personalizar la apariencia de la aplicación. La configuración debe conservarse en el
dispositivo.}
{Media}
{Sin relación directa}
{Cambiar el tema y el acento, recargar la aplicación y comprobar que ambos
valores se mantienen.}


\subsection{Requisitos no funcionales}
\setsecnumdepth{section}
\settocdepth{section}

Se recogen diez \textbf{requisitos no funcionales} generales. Para cada uno
se indica su alcance y las funciones relacionadas. Pese a no tener un campo
«Verificación», ya que en la mayoría de ellos no hay una metodología objetiva
para comprobar que están implementados de una forma correcta, todos ellos deben ser tenidos en cuenta durante la fase de \ac{QA}.

\requisitoNoFuncional
{RNF-1}
{Usabilidad y accesibilidad}
{\refObjetivo{OG-01}}
{Interfaz de usuario}
{La interfaz debe adaptarse a móvil y escritorio y cumplir el nivel AA de las
\ac{WCAG} 2.2~\cite{w3c_wcag_22}, incluyendo el uso con teclado, el foco visible
y la identificación de controles.
Los mensajes de error deben explicar el problema y cómo continuar.}
{Alta}
{\refRequisito{RF-1}, \refRequisito{RF-2}, \refRequisito{RF-3}, \refRequisito{RF-4} y
\refRequisito{RF-5}}


\requisitoNoFuncional
{RNF-2}
{Compatibilidad e internacionalización}
{\refObjetivo{OG-01}}
{Interfaz web}
{Las funciones deben conservar su comportamiento en los navegadores
admitidos. Los textos y mensajes de la interfaz deben estar disponibles de forma
coherente en español e inglés.}
{Media}
{\refRequisito{RF-1}, \refRequisito{RF-2}, \refRequisito{RF-3}, \refRequisito{RF-4} y
\refRequisito{RF-5}}

\requisitoNoFuncional
{RNF-3}
{Rendimiento y consumo de recursos}
{\refObjetivo{OG-01}, \refObjetivo{OG-02} y \refObjetivo{OG-03}}
{Toda la aplicación}
{Los tiempos de respuesta y el consumo de recursos deben mantenerse dentro
de los límites definidos para la carga prevista. Las operaciones más costosas
no deben bloquear el uso de otras funciones.}
{Alta}
{\refRequisito{RF-1}, \refRequisito{RF-2}, \refRequisito{RF-3},
\refRequisito{RF-4} y \refRequisito{RF-5}}

\requisitoNoFuncional
{RNF-4}
{Fiabilidad y recuperación}
{\refObjetivo{OG-02}, \refObjetivo{OG-03} y \refObjetivo{OG-04}}
{Datos persistentes y tareas}
{Los fallos de una tarea deben quedar registrados sin interrumpir trabajos
independientes ni detener toda la aplicación. Las tareas deben ser idempotentes,
de modo que puedan reintentarse sin generar duplicados ni aplicar varias
veces los mismos cambios.}
{Alta}
{\refRequisito{RF-1}, \refRequisito{RF-2}, \refRequisito{RF-3} y \refRequisito{RF-4}}

\requisitoNoFuncional
{RNF-5}
{Seguridad}
{\refObjetivo{OG-05}}
{Toda la aplicación}
{La aplicación debe proteger las cuentas y los datos frente a accesos y
modificaciones sin permiso. Los datos deben estar protegidos tanto al
guardarlos como al enviarlos. Esto incluye la protección de las sesiones de
usuario~\cite{owasp_session_management} y la prevención de ataques de
\ac{CSRF} en las peticiones que modifican datos~\cite{owasp_csrf}.}
{Alta}
{\refRequisito{RF-1}, \refRequisito{RF-2}, \refRequisito{RF-3} y \refRequisito{RF-4}}

\requisitoNoFuncional
{RNF-6}
{Privacidad y tratamiento de datos}
{\refObjetivo{OG-04} y \refObjetivo{OG-05}}
{Cuentas, manuales y conversaciones}
{El uso de los datos personales debe limitarse al mínimo necesario para el
uso del sistema.}
{Alta}
{\refRequisito{RF-1}, \refRequisito{RF-3} y \refRequisito{RF-4}}


\requisitoNoFuncional
{RNF-7}
{Calidad del procesamiento y las respuestas}
{\refObjetivo{OG-02} y \refObjetivo{OG-03}}
{Texto y respuestas generados}
{El texto extraído y corregido debe conservar las reglas del manual.
Por otro lado, se deben aplicar las medidas necesarias para evitar,
en la medida de lo posible, las alucinaciones del modelo \ac{LLM}~\cite{openai_hallucinations}.}
{Alta}
{\refRequisito{RF-3.4}, \refRequisito{RF-4.1}, \refRequisito{RF-4.3},
\refRequisito{RF-4.5} y \refRequisito{RF-4.6}}


\requisitoNoFuncional
{RNF-8}
{Mantenibilidad y pruebas}
{\refObjetivo{OG-01}}
{Código y proceso de integración}
{El código debe separar responsabilidades y disponer de pruebas de los
comportamientos principales y sus errores. Las comprobaciones del código y las
dependencias deben formar parte de la integración. La configuración y el
mantenimiento deben estar documentados.}
{Media}
{\refRequisito{RF-1}, \refRequisito{RF-2}, \refRequisito{RF-3}, \refRequisito{RF-4} y
\refRequisito{RF-5}}

\requisitoNoFuncional
{RNF-9}
{Portabilidad y reproducibilidad}
{\refObjetivo{OG-01}}
{Despliegue y ejecución de la aplicación}
{La aplicación debe poder desplegarse y ejecutarse en los entornos admitidos
mediante un procedimiento documentado.}
{Baja}
{\refRequisito{RF-1}, \refRequisito{RF-2}, \refRequisito{RF-3}, \refRequisito{RF-4} y
\refRequisito{RF-5}}

\requisitoNoFuncional
{RNF-10}
{Observabilidad}
{\refObjetivo{OG-01}, \refObjetivo{OG-02} y \refObjetivo{OG-03}}
{Supervisión y diagnóstico}
{La aplicación debe proporcionar información sobre el estado de sus servicios
y la ejecución de las tareas para facilitar la detección y el diagnóstico de
fallos. Los registros de error deben indicar cuándo se produjo el fallo, en qué
componente y durante qué operación o tarea.}
{Media}
{\refRequisito{RF-1}, \refRequisito{RF-2}, \refRequisito{RF-3},
\refRequisito{RF-4}, \refRequisito{RNF-5} y \refRequisito{RNF-6}}
